%%%%%%%%%%%%%%%%%%%%%%%%%%%%%%%%%%%%%%%%%%%%%
%%        math bold fonts                  %%
%%%%%%%%%%%%%%%%%%%%%%%%%%%%%%%%%%%%%%%%%%%%%
\newcommand{\bA}{\mathbf{A}}
\newcommand{\bB}{\mathbf{B}}
\newcommand{\bC}{\mathbf{C}}
\newcommand{\bD}{\mathbf{D}}
\newcommand{\bE}{\mathbf{E}}
\newcommand{\bF}{\mathbf{F}}
\newcommand{\bG}{\mathbf{G}}
\newcommand{\bH}{\mathbf{H}}
\newcommand{\bI}{\mathbf{I}}
\newcommand{\bJ}{\mathbf{J}}
\newcommand{\bK}{\mathbf{K}}
\newcommand{\bL}{\mathbf{L}}
\newcommand{\bM}{\mathbf{M}}
\newcommand{\bN}{\mathbf{N}}
\newcommand{\bO}{\mathbf{O}}
\newcommand{\bP}{\mathbf{P}}
\newcommand{\bQ}{\mathbf{Q}}
\newcommand{\bR}{\mathbf{R}}
\newcommand{\bS}{\mathbf{S}}
\newcommand{\bT}{\mathbf{T}}
\newcommand{\bU}{\mathbf{U}}
\newcommand{\bV}{\mathbf{V}}
\newcommand{\bW}{\mathbf{W}}
\newcommand{\bX}{\mathbf{X}}
\newcommand{\bY}{\mathbf{Y}}
\newcommand{\bZ}{\mathbf{Z}}

\newcommand{\ba}{\mathbf{a}}
\newcommand{\bb}{\mathbf{b}}
\newcommand{\bc}{\mathbf{c}}
\newcommand{\bd}{\mathbf{d}}
\newcommand{\be}{\mathbf{e}}
\newcommand{\bbf}{\mathbf{f}}
\newcommand{\bg}{\mathbf{g}}
\newcommand{\bh}{\mathbf{h}}
\newcommand{\bi}{\mathbf{i}}
\newcommand{\bj}{\mathbf{j}}
\newcommand{\bk}{\mathbf{k}}
\newcommand{\bl}{\mathbf{l}}
\newcommand{\bm}{\mathbf{m}}
\newcommand{\bn}{\mathbf{n}}
\newcommand{\bo}{\mathbf{o}}
\newcommand{\bp}{\mathbf{p}}
\newcommand{\bq}{\mathbf{q}}
\newcommand{\br}{\mathbf{r}}
\newcommand{\bs}{\mathbf{s}}
\newcommand{\by}{\mathbf{y}}
\newcommand{\bu}{\mathbf{u}}
\newcommand{\bv}{\mathbf{v}}
\newcommand{\bw}{\mathbf{w}}
\newcommand{\bx}{\mathbf{x}}
\newcommand{\bby}{\mathbf{y}}
\newcommand{\bz}{\mathbf{z}}


%%%%%%%%%%%%%%%%%%%%%%%%%%%%%%%%%%%%%%%%%%%%%
%%        math cal fonts                  %%
%%%%%%%%%%%%%%%%%%%%%%%%%%%%%%%%%%%%%%%%%%%%%
\newcommand{\cA}{\mathcal{A}}
\newcommand{\cB}{\mathcal{B}}
\newcommand{\cC}{\mathcal{C}}
\newcommand{\cD}{\mathcal{D}}
\newcommand{\cE}{\mathcal{E}}
\newcommand{\cF}{\mathcal{F}}
\newcommand{\cG}{\mathcal{G}}
\newcommand{\cH}{\mathcal{H}}
\newcommand{\cI}{\mathcal{I}}
\newcommand{\cJ}{\mathcal{J}}
\newcommand{\cK}{\mathcal{K}}
\newcommand{\cL}{\mathcal{L}}
\newcommand{\cM}{\mathcal{M}}
\newcommand{\cN}{\mathcal{N}}
\newcommand{\cO}{\mathcal{O}}
\newcommand{\cP}{\mathcal{P}}
\newcommand{\cQ}{\mathcal{Q}}
\newcommand{\cR}{\mathcal{R}}
\newcommand{\cS}{\mathcal{S}}
\newcommand{\cT}{\mathcal{T}}
\newcommand{\cU}{\mathcal{U}}
\newcommand{\cV}{\mathcal{V}}
\newcommand{\cW}{\mathcal{W}}
\newcommand{\cX}{\mathcal{X}}
\newcommand{\cY}{\mathcal{Y}}
\newcommand{\cZ}{\mathcal{Z}}



% \inmath{material}
%    Make sure that {material} is inside math mode.
\def\inmath#1{\relax\ifmmode#1\else$#1$\fi}

% \@<token>
%    Will print the argument in \cal.
\def\@#1{\inmath{{\cal #1}}}

% \operators{<list>}
% \variables{<list>}
%    A convenient way of defining a set of variables or functions for
%    use in math. <list> consists of a \\ separated list of pairs
%    <token><text> where <token> is the token that prints <text> as an
%    operator/variable in math-mode (using \mathop or \mathord).
%
%    See \det and \gcd below for examples.
%
\outer\def\operators#1{\begingroup
        \def\\##1##2{\gdef##1{{\mathop{\rm ##2}}}}\relax
        \\#1\endgroup
}

\outer\def\variables#1{\begingroup
        \def\\##1##2{\gdef##1{{\mathord{\it ##2}}}}\relax
        \\#1\endgroup
}

% \ntimes#1#2
% make #1 repetitions of #2
\def\ntimes#1#2{{\count255=1\loop#2\ifnum\count255<#1\advance\count255 by1\repeat}}

% \boxit{material}
%    This one boxes {material} using rules and is a vbox itself
\long\def\boxit#1{\vbox{\hrule\hbox{\vrule\kern3pt\vbox{\kern3pt#1\kern3pt}\kern3pt\vrule}\hrule}}

\long\def\entry#1#2\par{\goodbreak\item{$\bullet$}\strut{#1}\par\nobreak{\narrower\vbox{#2\strut}}\par}

% \th{1}
%    Adds apropriate ``th'' extension to digits (``1st'', ``2nd'', ``3d'', ....)
\def\th#1{#1\ifcase#1 th\or st\or nd\or d\else th\fi}


% \smiley
% \frowny
%    Here we have a smile face (and a frowny face) to use together with
%    text in horisontal mode
\begingroup
        \def\facewith#1{$\bigcirc\mskip-13.3mu{}^{..}
        \mskip-11mu\scriptscriptstyle#1\ $}
        \xdef\frowny{\facewith\frown}
        \xdef\smiley{\facewith\smile}
\endgroup


% \underrightarrow{material}
% \underleftarrow{material}
%    Insert a right or left arrow under {material}, does not affect the baseline
\def\underrightarrow#1{\vtop{\ialign{##\crcr
        \hbox{#1}\crcr
        \noalign{\kern-1pt\nointerlineskip}
        \rightarrowfill\crcr}}}
\def\underleftarrow#1{\vtop{\ialign{##\crcr
        \hbox{#1}\crcr
        \noalign{\kern-1pt\nointerlineskip}
        \leftarrowfill\crcr}}}


% \addtobox\foo{bar}
%    This one adds {bar} to the box represented by \foo.
%    If the box is void, the material is inserted into the box
%
%    NOTE! The box \foo has to be of the same type as the {bar} material
\def\addtobox#1#2{\setbox#1=\ifvoid#1
  {#2}
  \else\ifvbox#1
    \vbox{\unvbox#1\relax#2}
    \else\ifhbox#1
      \hbox{\unhbox#1\relax#2}
    \fi\fi\fi}

% \today
%    This one constructs todays date in the format ``November 19, 1999''
\def\today{\ifcase\month\or January\or February\or March\or April\or May\or June\or July\or
        August\or September\or October\or November\or December\fi
        \space\number\day, \number\year}

% \idag
%    A swedish version of the \today macro.
\def\idag{\number\day\space
        \ifcase\month\or Januari\or Februari\or Mars\or April\or Maj\or Juni\or Juli\or
        Augusti\or September\or Oktober\or November\or December\fi
        \space\number\year}


% \mwidthof{foo}
% \mheightof{foo}
%    Prints the height or the width of {foo} as a message
\def\mwidthof#1{\setbox0=\hbox{#1}\message{\the\wd0}}
\def\mheightof#1{\setbox0=\hbox{#1}\message{\the\ht0}}


% \widthof{foo}
% \heightof{foo}
%    Expands to the width or the height of {foo}.
\def\widthof#1{\setbox0=\hbox{#1}\the\wd0}
\def\heightof#1{\setbox0=\hbox{#1}\the\ht0}


% \nextpar{sequence}
%    Will enter {sequence} at the beginning of the next paragraph.
\def\nextpar#1{\everypar={#1\everypar={}}}


% \dropnextpar
%    Will enlarge the first letter of the next paragraph and drop it
%    into the following text in the manner that some storybooks (and
%    old bibles) do
\def\dropnextpar{\begingroup\parindent=0pt
        \font\bigfont=cmr10 at 2\baselineskip
        \everypar={\futurelet\a\cont}}
\def\cont#1#2{\setbox0=\hbox{\bigfont#1\kern1.5pt}\relax
        \setbox1=\hbox{#2}\hangindent=\wd0
        \hangafter=-\ht0\advance\hangafter by -\baselineskip
        \divide\hangafter by\baselineskip
        \noindent\llap{\vbox to\the\ht1{\vskip-.5pt\box0\vss}}#2\endgroup}

% \registered
%    Registered trademark symbols
\def\registered{\raise.5ex\hbox{\ooalign{\hfil\raise.06ex
        \hbox{$\scriptstyle\rm R$}\hfil\crcr\mathhexbox20D}}}


%%% Local Variables: 
%%% mode: latex
%%% TeX-master: t
%%% End: 
